% القسم: المناهج
% الفصل: الاستقراء
% المبحث: التعريفات الاستقرائية

\documentclass[../../../include/open-logic-section]{subfiles}

\begin{document}

\olfileid{mth}{ind}{idf}

\olsection{التعريفات الاستقرائية}

كثير من أصناف الأشياء في المنطق يُعرّف \emph{استقرائيًّا}: توضع قواعد
تعيّن أفراد الصنف الأولين، وتبين كيف تُستحدث أفراد منه من أفراد سبق
حصولها. ومن ذلك تعريف ضروب من متواليات الرموز، كحدود لغة و!!{formula}
فيها. ولنأخذ مثالًا يسيرًا: سلاسل تتألف من الحروف $\mathrm{a}$ و
$\mathrm{b}$ و$\mathrm{c}$ و$\mathrm{d}$، ومن الرمز~$\circ$،
والقوسين $[$ و$]$، كالسلاسل
«$[[\mathrm{c} \circ \mathrm{d}][$» و«$[\mathrm{a}[]\circ]$»
و«$\mathrm{a}$» و«$[[\mathrm{a} \circ \mathrm{b}]\circ
\mathrm{d}]$». وفي الأوليين خلل ظاهر: أقواس الأولى غير «متوازنة»،
والرمز «$\circ$» في الثانية لا «يصل» بين عبارتين تامتين في نفسيهما.
وأما الثالثة والرابعة فأقوم؛ فكل «$[$» يقابله قوس إغلاق «$]$» إذا
وجدت الأقواس، وكل $\circ$ واقع بين عبارتين «سليمتين» قد أحاط بهما زوج
من الأقواس.

ونريد ضبط معنى «الحد السليم». فالحرف المنفرد سليم ابتداءً. وما جاوز
الحرف الواحد فلا بد أن يكون على صورة «$[t \circ s]$»، على أن يكون
$s$ و$t$ سليمين كل واحد في نفسه. وبالعكس، إذا سلم $t$ و$s$، أمكن
إنشاء حد سليم جديد بوضع $\circ$ بينهما وإحاطتهما بقوسين. وبهاتين
العمليتين \emph{نعرّف} مجموعة الحدود السليمة؛ فهذا هو
\emph{التعريف الاستقرائي}.

\begin{defn}[الحدود السليمة]
  تعرّف مجموعة \emph{الحدود السليمة} استقرائياً كما يأتي:
  \begin{enumerate}
  \item كل حرف من $\mathrm{a}$ و$\mathrm{b}$ و$\mathrm{c}$
    و$\mathrm{d}$ حد سليم.
  \item إذا كان $s_1$ و$s_2$ حدين سليمين، فإن
    $[s_1 \circ s_2]$ حد سليم أيضاً.
  \item لا شيء آخر حد سليم.
  \end{enumerate}
\end{defn}

ومؤدى التعريف أن الشيء حد سليم إذا وفقط إذا أمكن بناؤه بالشرطين (1)
و(2) في خطوات متناهية. ففي الطبقة الأولى نحصل الحدود المؤلفة من حرف
واحد، وهي:
\[
\mathrm{a}, \mathrm{b}, \mathrm{c}, \mathrm{d}
\]
وفي الطبقة الثانية نعمل (2) في الحدود السابقة، فنحصل:
\[
[\mathrm{a} \circ \mathrm{a}], [\mathrm{a} \circ \mathrm{b}],
[\mathrm{b} \circ \mathrm{a}], \dots, [\mathrm{d} \circ \mathrm{d}]
\]
وهذه جميع تراكيب الحرفين. ثم نعود في الطبقة الثالثة إلى (2)، فنعمله
في أي حدين سليمين حصلا قبلها. فتظهر حدود جديدة، منها
$[\mathrm{a} \circ [\mathrm{a} \circ \mathrm{a}]]$---حيث $t$ هو
$\mathrm{a}$ من الخطوة~1 و$s$ هو
$[\mathrm{a} \circ \mathrm{a}]$ من الخطوة~2---ومثل
$[[\mathrm{b} \circ \mathrm{c}] \circ [\mathrm{d} \circ
\mathrm{b}]]$، المؤلف من الحدين $[\mathrm{b} \circ \mathrm{c}]$
و$[\mathrm{d} \circ \mathrm{b}]$ من الطبقة~2. وهكذا يتوالى البناء،
ويمنع البند (3) دخول ما لم يُبن بهذه السبيل في مجموعة الحدود السليمة.

ولم نبرهن بعد أن كل سلسلة «تبدو» سليمة تدخل في التعريف. غير أن العكس
ظاهر: فكل ما يبنيه التعريف «يبدو سليمًا»؛ لأن أقواسه متوازنة، ولأن
$\circ$ يصل بين جزأين سليمين في نفسيهما.

وأعظم منفعة للتعريف الاستقرائي أنه يحصر، عند إثبات حكم لجميع الحدود
السليمة، الحالات التي يجب استيفاؤها. فإذا كان $t$ حدًّا سليمًا، حصر
التعريف صورة $t$: إما حرف منفرد، وإما $[s_1 \circ s_2]$ لحدين سليمين
$s_1$ و~$s_2$. وبالبند (3) يمتنع كل احتمال سواهما.

وللاستقراء القوي شأن خاص في إثبات القضايا على جميع أفراد مجموعة
معرّفة استقرائيًّا. فلو أردنا، مثلًا، بيان أن عدد الرموز $[$ في كل حد
سليم طوله~$n$ هو~$< n/2$، أمكن أن نجعله قضية في جميع~$n$: لكل $n$،
عدد الرموز $[$ في أي حد سليم طوله~$n$ هو $< n/2$.

\begin{prop}
  لكل $n$، عدد الرموز $[$ في حد سليم طوله~$n$ هو $< n/2$.
\end{prop}

\begin{proof}
طريق الاستقراء القوي أن نثبت الشرطية الآتية:
\begin{quote}
  إذا كان كل حد سليم طوله~$l$ لا يحتوي من الرموز $[$ إلا عددًا أقل من
  $l/2$ لكل $l < k$، فإن كل حد سليم طوله~$k$ لا يحتوي منها إلا
  عددًا أقل من $k/2$.
\end{quote}
ولإثبات الشرطية نفرض مقدمها: أن الحدود السليمة ذات الطول~$l$ فيها من
الرموز $[$ عدد أقل من $l/2$، لأي $l<k$. وهذا هو فرض الاستقراء؛
والمطلوب نقل الحكم إلى الحدود السليمة ذات الطول~$k$.

فليكن $t$ حدًّا سليمًا طوله~$k$. ولما كان تعريف الحدود استقرائيًّا،
انحصرت صورة $t$ في حالتين: (1)~أن يكون حرفًا منفردًا، أو (2)~أن يكون
  أن يكون~$t$ على صورة $[s_1 \circ s_2]$ لحدين سليمين $s_1$ و~$s_2$.
\begin{enumerate}
\item $t$ حرف. فحينئذ $k = 1$، وعدد الرموز $[$ في $t$ هو~$0$.
ولما كان $0 < 1/2$ ثبت الحكم.
\item $t$ هو $[s_1 \circ s_2]$ لبعض حدين سليمين $s_1$ و~$s_2$.
  ليكن $l_1$ طول~$s_1$ و$l_2$ طول~$s_2$. عندئذ يكون طول~$t$، وهو
  $k$، مساوياً $l_1+l_2+3$، أي مجموع طولي $s_1$ و$s_2$ مع الرموز
  الثلاثة $[$ و$\circ$ و$]$. ولما كان $l_1+l_2+3$ أكبر دائماً من
  $l_1$، فإن $l_1 < k$. وبالمثل $l_2 < k$. وهذا يعني أن فرض
  الاستقراء ينطبق على الحدين $s_1$ و~$s_2$: فعدد الرموز $[$ في
  $s_1$، وليكن~$m_1$، أقل من $l_1/2$، وعددها في~$s_2$، وليكن~$m_2$،
  أقل من $l_2/2$.

  عدد الرموز $[$ في $t$ هو عددها في~$s_1$ مضافاً إليه عددها في~$s_2$
  ثم~$1$، أي إنه $m_1 + m_2 + 1$. وبما أن $m_1 < l_1/2$ و
  $m_2 < l_2/2$، فلدينا:
  \[
  m_1 + m_2 + 1 < \frac{l_1}{2} + \frac{l_2}{2} + 1 = \frac{l_1+l_2+2}{2} < \frac{l_1+l_2+3}{2} = k/2.
  \]
\end{enumerate}
وقد ظهر في كلتا الحالتين، اعتمادًا على فرض الاستقراء، أن عدد الرموز
$[$ في $t$ أقل من $k/2$؛ فتثبت القضية بالاستقراء القوي.
\end{proof}

\begin{prob}
  لتكن مجموعة الحدود فائقة السلامة معرّفة كما يأتي:
  \begin{enumerate}
  \item كل حرف من $\mathrm{a}$ و$\mathrm{b}$ و$\mathrm{c}$
    و$\mathrm{d}$ حد فائق السلامة.
  \item إذا كان $s$ حداً فائق السلامة، فإن $[s]$ كذلك.
  \item إذا كان $s_1$ و$s_2$ حدين فائقي السلامة، فإن
    $[s_1 \circ s_2]$ كذلك.
  \item لا شيء آخر حد فائق السلامة.
  \end{enumerate}
  أثبت أن عدد الرموز $[$ في حد فائق السلامة~$t$ طوله~$n$ لا يزيد على
  $n/2 +1$.
\end{prob}

\end{document}
